% !TEX root = ../../cookbook.tex
\newpage
\recipe{Pizza napoletana}{Dalla \href{https://www.youtube.com/watch?v=xqVgJHThsco&t=195s}{ricetta di Brian Lagerstrom}. Per 4 pizze.}{}{} \label{pizza_napolitana}

\subsection*{Ingredienti}
\begin{ingredients}
	\item 450 g di acqua
	\item 6 pizzichi di lievito in polvere (\sfrac{1}{8} di cucchiaino)
	\item 665 g di farina 00
	\item 20 g di sale
\end{ingredients}


\medskip

% --- Preparation ---
Mettere il \textbf{lievito} nell'\textbf{acqua}\footnote{Per attivare bene il lievito, metterlo nell'acqua calda con un cucchiaino di zucchero o farina per 5-10 minuti, finché schiuma.}, poi aggiungere \textbf{farina} e \textbf{sale}. 
Mischiare bene con un cucchiaio, poi impastare con una planetaria\footnote{In assenza di una planetaria, si può impastare anche con le mani, ma come per il pane, il risultato sarà una crosta meno alveolata e soffice, e in generale meno professionale.} finché ben uniforme ed elastica. 
Coprire e lasciar riposare 30 minuti. 

Eseguire qualche piega di rinforzo per sviluppare il glutine\footnote{\textit{Pieghe di rinforzo}: senza togliere l'impasto dal suo contenitore, pizzica un lato del panetto e tiralo verso il lato opposto. Ripeti sui 4 lati del panetto, formando una `X'. Gira il panetto sulla chiusura prima di lasciarlo riposare di nuovo. Brian Lagerstrom ha molti video, incluso quello nel link in questa ricetta, sotto al titolo, dove mostra come si esegue una piega di rinforzo.\label{pieghe_rinforzo}}, poi coprire e lasciar riposare altri 30 minuti.

Eseguire un'altra piega di rinforzo, poi coprire e lasciar lievitare fuori dal frigo per 12h.

Separare il panetto in 4 palle, e metterle in frigo in contenitori oleati per almeno 2h, ma fino a 5 giorni.

Rimuovere dal frigo 3h prima di fare la pizza.
Spianare bene le pizze su della farina, guarnire a piacere\footnote{Suggerisco, dopo aver spianato la pizza (rigorosamente con le mani), di spostarla e guarnirla rapidamente su una pala di alluminio o legno ben infarinata. Il prelevamento della pizza già adornata direttamente dal tavolo richiede non poca destrezza e, in mani inesperte, spesso risulta nella distruzione della stessa. Inoltre, cerca di lasciare la pizza sulla pala il meno tempo possibile, soprattutto dopo aver aggiunto ingredienti bagnati (come il pomodoro), poichè l'umidità farà che la pizza appiccichi alla pala.}\footnote{In caso che il lettore, o uno degli aspiranti attavolati, preveda di guarnire la pizza con dell'\textit{ananas}, voglio sottolineare che questo è un atto assolutamente barbarico, meritevole di publica umiliazione.}, poi cuocere in forno preriscaldato, idealmente su una pietra da forno, per circa 10 minuti (il sotto deve iniziare ad avere delle zone belle scure o lievemente sbruciacchiate). 

% --- Description ---
\begin{recipeblock}
\noindent \textit{``A Napoli hanno inventato la pizza. \\ 
E tutti suonano il clacson, \\ perchè vogliono avere la pizza.'' \\ \normalfont \hspace{1.5cm} \textemdash i fratelli Rizzello.}
\end{recipeblock}
\vfill
